\documentclass[10pt,a4paper]{amsart}
\usepackage[margin=19mm]{geometry}
\usepackage{amsmath,amssymb,mathtools}
\usepackage[scr=rsfs,cal=euler]{mathalfa}
\usepackage{xfrac}
\usepackage[unicode,hidelinks,bookmarks=false]{hyperref}
\newcommand{\ie}{i.e.\ }
\newcommand{\cf}{cf.\ }
\newcommand{\resp}{respectively\ }
\newcommand{\Subsection}{\subsection}
\providecommand{\nobreakdash}{\nobreak\hbox{-}}
\setlength{\emergencystretch}{2em}
\multlinegap0pt
\usepackage{bm}
\usepackage{bbm}
\usepackage{dsfont}
\usepackage{xspace}
\usepackage{cancel}
\usepackage{mathtools}
\usepackage{amsthm}
\makeatletter
\@ifundefined{lemma}{\newtheorem{lemma}{Lemma}}{}
\makeatother
\title[On Tent Spaces for the Gaussian Measure]{On Tent Spaces for the Gaussian Measure}
\author{Liliana Forzani, Roberto Scotto,, Wilfredo Urbina}
\date{Source: arXiv:2509.16148v2 (CC BY 4.0)}
\begin{document}
\maketitle

\noindent\emph{Source and licence.} On Tent Spaces for the Gaussian Measure. Version arXiv:2509.16148v2 (CC BY 4.0), distributed under the Creative Commons Attribution 4.0 International licence, as recorded in the arXiv licence metadata for this exact version and shown on the abstract page https://arxiv.org/abs/2509.16148v2. Full rights notice: licenses/arxiv-2509.16148-CC-BY-4.0.txt.\par\smallskip

\noindent\textit{Editorial scope.} Counted unit: the Lemma whose source label is \texttt{lema: desigualdad integral, relacion entre eta y beta} in the article (source0.tex lines 749--755, source label \texttt{lema: desigualdad integral, relacion entre eta y beta}), delivered together with the article's own complete original proof of it (source0.tex lines 757--802, one \texttt{proof} environment of 46 lines). No mathematics is written, repaired or re-derived: statement and proof are the article's own text line for line. Required context retained uncounted: comparacion (lines 550--553). Every definition, display and label of those lines is retained unchanged. Omitted from this package and not redistributed: 1--549, 554--748, 756--756, 803--1474. Nothing inside a retained excerpt is omitted or altered: every retained body is delivered exactly as the article's lines are written, so no omission receipt of any kind accompanies this package. Citations: the retained excerpts carry no citation command at all, so no bibliography entry of the article is transcribed and no reference is redistributed. Reference adaptation: none was needed and none was made; every reference inside the delivered unit resolves inside it, so no unresolved reference or citation is produced. Printed slips kept literal: no printed word, symbol, hypothesis or number of the counted statement or of its proof was altered. Layout and class: the article's own class is \texttt{amsart} with packages including amsfonts, amssymb, amsmath, amsthm, hyperref, color, float, enumerate, enumitem, mathrsfs, url, graphicx, footmisc, esint; the wrapper is the lane's pinned \texttt{amsart} editorial wrapper at 10pt on A4, which restates the theorem-like environments the retained text needs and adds the article's own macro definitions that the retained text needs (none), with extra wrapper packages (bm, bbm, dsfont, xspace, cancel, mathtools). The delivered numbering is the unit's own. Assets: no retained span contains a figure environment, a graphics-inclusion command, a PDF-page inclusion, a table or a TikZ picture; no image file is copied into this package, the package declares no asset dependency and no image file is redistributed with it. Licence: version arXiv:2509.16148v2 of this article is distributed under the Creative Commons Attribution 4.0 International licence, as recorded in the arXiv licence metadata for this exact version and shown on the abstract page https://arxiv.org/abs/2509.16148v2. A full rights notice is delivered at licenses/arxiv-2509.16148-CC-BY-4.0.txt. Characters: every retained excerpt is pure ASCII, so the delivered engine, XeLaTeX with its default Latin Modern text font, reports no missing character.
  \begin{lemma}\label{comparacion}
Let $b>0$ be given. If $|x-y|<b\, m(y)$ then 
$$m(y)<(b+1)m(x), \quad  \text{and} \quad m(x)<(b+1)m(y).$$
\end{lemma}

\begin{lemma}\label{lema: desigualdad integral, relacion entre eta y beta}
For every $\eta\in (0,1)$ there exists $\bar{\eta}\in (0,1)$ such that for every measurable subset $A$ of $\mathbb R^n$ and every non-negative measurable function $H$ on $\mathbb R^{n+1}_+$, we have
\begin{align*}
    \iint_{R^{(1-\eta)\alpha, (1-\eta)\beta}_\gamma\left(A_{\beta(1+\beta)}^{[\bar{\eta]}}\right)} & H(y,t) d\gamma(y) \frac{dt}{t} \\
    &\lesssim \int_A \left(\iint_{\Gamma_\gamma^{\alpha, \beta}(x)} \frac{H(y,t)}{\gamma(B(y,\alpha t\wedge m_\beta(y))} d\gamma(y) \frac{dt}{t}\right) d\gamma(x).
\end{align*}
\end{lemma}

\begin{proof}
By applying Fubini's Theorem, we get
\begin{align*}
    \int_A &\left(\iint_{\Gamma_\gamma^{\alpha, \beta}(x)} \frac{H(y,t)}{\gamma(B(y,\alpha t\wedge m_\beta(y)))} d\gamma(y) \frac{dt}{t}\right) d\gamma(x) \\
    &= \iint_{\mathbb R^{n+1}_+} \left(\int_A \mathcal{X}_{B(y,\alpha t\wedge m_\beta (y))}(x) d\gamma(x)\right)  \frac{H(y,t)}{\gamma(B(y,\alpha t\wedge m_\beta(y)))} d\gamma(y) \frac{dt}{t}\\
    &= \iint_{\mathbb R^{n+1}_+} \frac{\gamma(A\cap B(y,\alpha t\wedge m_\beta (y)))}{\gamma(B(y,\alpha t\wedge m_\beta (y)))} H(y,t) d\gamma(y) \frac{dt}{t}.
\end{align*}
We wish to find a suitable subset of ${\mathbb R}^{n+1}_+$ so that the quotient $\frac{\gamma(A\cap B(y,\alpha t\wedge m_\beta(y)))}{\gamma(B(y,\alpha t\wedge m_\beta(y)))}$ be bounded below by a constant on such a subset and  in this way to obtain the desired inequality.

First, notice that for every $(y,t)\in R_{(1-\eta)\alpha, (1-\eta)\beta}\left(A^{[\bar{\eta}]}_{\beta (1+\beta)}\right)$, there exists $x\in A^{[\bar{\eta}]}_{\beta(1+\beta)}$ such that ${|x-y|<(1-\eta)(\alpha t\wedge m_\beta(y))}$. Thus, from Lemma~\ref{comparacion} with $b=\beta$, $ m(y)\leq (1+\beta)\, m(x)$. 

From the above fact, $\alpha t\wedge m_\beta(y) \leq \beta\,  m(y)\leq \beta (1+\beta)\,  m(x)$, so we have $B(x,\alpha t\wedge m_\beta(y))\in \mathscr{B}_{\beta (1+\beta)}$, and since $x\in A^{[\bar{\eta}]}_{\beta (1+\beta)}$, then 
\[\frac{\gamma(A\cap B(x,\alpha t\wedge m_\beta(y)))}{\gamma(B(x,\alpha t\wedge m_\beta(y)))}\geq \bar{\eta}.\]

Therefore,
\begin{align*}
  &\gamma(A\cap B(y,\alpha t\wedge m_\beta(y)))\geq \gamma(A\cap B(x,\alpha t\wedge m_\beta(y)))\\ &\ \ \ \ -\gamma(B(x,\alpha t\wedge m_\beta(y))\cap B^c(y,\alpha t\wedge m_\beta(y)))\\
  &\geq \bar{\eta} \gamma(B(x,\alpha t\wedge m_\beta(y)))-\gamma(B(x,\alpha t\wedge m_\beta(y))\cap B^c(y,\alpha t\wedge m_\beta(y))).
\end{align*}

From the relationship between $x$ and $y$, we get $B(x,\eta (\alpha t\wedge m_\beta(y)))\subseteq B(y,\alpha t\wedge m_\beta(y))$. On the other hand, from the doubling property of the Gaussian measure on admissible balls from $\mathscr{B}_{\beta (1+\beta)}$, we know that there exists a constant $C=C(1/\eta, \beta)>1$ such that
\[\gamma(B(x,\alpha t\wedge m_\beta(y)))\leq C \gamma(B(x,\eta (\alpha  t\wedge m_\beta(y)))).\]
Hence,
\begin{align*}
\gamma(B(x,\alpha t\wedge m_\beta(y))\cap B(y,\alpha t\wedge m_\beta(y)))&\geq \gamma(B(x,\eta (\alpha t\wedge m_\beta(y))))\\ &\geq \frac{1}{C}\gamma(B(x, \alpha t\wedge m_\beta(y))).\end{align*}
Consequently, 
\begin{align}\label{eq: gamma de F intersec la bola}
  \nonumber \gamma(A\cap B(y,\alpha t\wedge m_\beta(y)))&\geq \bar{\eta} \gamma(B(x,\alpha t\wedge m_\beta(y))) \nonumber \\ &\ \ \ \ -\left(1-\frac{1}{C}\right)\gamma(B(x,\alpha t\wedge m_\beta(y)))\nonumber\\
  &=\left(\bar{\eta}-1+\frac{1}{C}\right)\gamma(B(x,\alpha t\wedge m_\beta(y))).
\end{align}

Again, since $|y-x|<\beta\, m(y)$ and $ m(y)\leq (1+\beta) m(x)$,
\begin{align*}
    \left||y|^2-|x|^2\right|&=\left||y|-|x|\right|\left(|x|+|y|\right)\leq \beta\, m(y)(|x|+|y|) \\
    & \leq \beta \, (1+\beta) m(x)|x|+ \beta \, m(y)|y|\\
    &\leq \beta\, (2+\beta).
\end{align*}
Then, there exists a positive constant $K=e^{-3\beta(2+\beta)}$ such that $\gamma(B(x,\alpha t\wedge m_\beta(y)))\geq K \gamma(B(y,\alpha t\wedge m_\beta(y)))$. Using this inequality in \eqref{eq: gamma de F intersec la bola}, we get 
\[\gamma(A\cap B(y,\alpha t\wedge m_\beta(y)))\geq \left(\bar{\eta}-1+\frac{1}{C}\right)K\, \gamma(B(y,\alpha t\wedge m_\beta(y))). \] Let us call $\lambda=\left(\bar{\eta}-1+\frac{1}{C}\right)K,$ then 
by choosing $\bar{\eta}$ such that $1-\frac{1}{C}<\bar{\eta} <1$, we get the desired estimate
\[\frac{\gamma(A\cap B(y,\alpha t\wedge m_\beta(y)))}{\gamma(B(y,\alpha t\wedge m_\beta(y)))}\geq \lambda,\]
for every $(y,t)\in R_{(1-\eta)\alpha, (1-\eta)\beta}\left(A^{[\bar{\eta}]}_{\beta(1+\beta)}\right)$. Therefore, 
\begin{align*} &\iint_{R_{(1-\eta)\alpha, (1-\eta)\beta}\left(A^{[\bar{\eta}]}_{\beta(1+\beta)}\right)} H(y,t) d\gamma(y) \frac{dt}{t}    \\
& \;\;\; \leq \frac{1}{\lambda} \int_A \left(\iint_{\Gamma_\gamma^{\alpha, \beta}(x)} \frac{H(y,t)}{\gamma(B(y,\alpha t\wedge m_\beta(y))} d\gamma(y) \frac{dt}{t}\right) d\gamma(x).\qedhere
\end{align*}
\end{proof}

\end{document}
